\begin{definition}{suite de Cauchy}{}
    Soit $(E, ||\cdot||)$ un evn, et soit $\alpha_n$ une suite de $E$. 

    On dit que $\alpha_n$ est une \emph{suite de Cauchy} si :
    \[
        \forall \varepsilon > 0, \exists N \in \mathbb{N}, \forall m,n \geqslant N, ||\alpha_n - \alpha_m|| < \varepsilon
    \]
\end{definition}

\begin{propriete}{}{}
    \begin{itemize}
        \item Toute suite convergente est une suite de Cauchy.
        \item Toute suite de Cauchy est bornée.
        \item Si une suite de Cauchy possède une valeur d'adhérence, alors elle converge
    \end{itemize}
\end{propriete}

\begin{theoreme}{convergence suites de Cauchy en dimension finie}{}
    Toute suite de Cauchy dans un evn de dimension finie converge.
\end{theoreme}

\begin{proof}
    Soit $(x_n)$ de Cauchy, alors $(x_n)$ est bornée, 
    
    donc par Bolzano-Weierstrass, elle admet une valeur d'adhérence, 
    donc elle converge.
\end{proof}

\begin{definition}{evn complet}{}
    Un evn $(E, ||\cdot||)$ est dit \emph{complet} si toute suite de Cauchy de $E$ converge dans $E$.
\end{definition}

\begin{remarque}
    Pour montrer qu'une suite diverge, toujours passer par l'absurde 
\end{remarque}

\subsubsection{Montrer qu'un evn est complet}

Soit $(E, ||\cdot||)$ un evn, et soit $(x_n)$ une suite de Cauchy de $E$.

\begin{enumerate}
    \item Construire un $x$ candidat à la limite de $(x_n)$
    \item Montrer que $x_n \in E$
    \item Montrer que $||x_n - x|| \to 0$
\end{enumerate}

\begin{proposition}{inégalité de Hölder}{}
    Soit $p, q \geqslant 1$ tels que $\frac{1}{p} + \frac{1}{q} = 1$. Alors $\forall f \in L^p(a,b)$ et $\forall g \in L^q(a,b)$, 
    \[
        ||f(t) g(t)||_1 dt \leqslant ||f||_p ||g||_q
    \]
    C'est à dire 
    \[
        \int_a^b |f(t) g(t)| dt \leqslant \left( \int_a^b |f(t)|^p dt \right)^{\frac{1}{p}} \left( \int_a^b |g(t)|^q dt \right)^{\frac{1}{q}}
    \]
\end{proposition}

\begin{proof}
    Avec inégalité de Jensen. 
\end{proof}

% \restaurergeometrie

\begin{proposition}{inégalité de Minkovski}{}
    Soit $p \geqslant 1$. Alors $\forall f, g \in L^p(a,b)$, 
    \[
        ||f + g||_p \leqslant ||f||_p + ||g||_p
    \]
    C'est à dire 
    \[
        \left( \int_a^b |f(t) + g(t)|^p dt \right)^{\frac{1}{p}} \leqslant \left( \int_a^b |f(t)|^p dt \right)^{\frac{1}{p}} + \left( \int_a^b |g(t)|^p dt \right)^{\frac{1}{p}}
    \]

\end{proposition}

\begin{proof}
    Inéglité triangulaire pour les normes et inégalité de Hölder.

    $${\displaystyle {
    \begin{aligned}\|f+g\|_{p}^{p}&=\int |f+g|^{p}\mathrm {d} \mu \\&\leq \int (|f|+|g|)|f+g|^{p-1}\mathrm {d} \mu \\&=\int |f||f+g|^{p-1}\mathrm {d} \mu +\int |g||f+g|^{p-1}\mathrm {d} \mu \\&\leq \left(\left(\int |f|^{p}\mathrm {d} \mu \right)^{1/p}+\left(\int |g|^{p}\mathrm {d} \mu \right)^{1/p}\right)\left(\int |f+g|^{(p-1)\left({\frac {p}{p-1}}\right)}\mathrm {d} \mu \right)^{1-{\frac {1}{p}}}\\&=(\|f\|_{p}+\|g\|_{p}){\frac {\|f+g\|_{p}^{p}}{\|f+g\|_{p}}},
    \end{aligned}}}$$
\end{proof}

\begin{proposition}{}{}
    $\forall p \geqslant 1$, $(L^p(\N), ||\cdot||_p)$ est un evn complet.
\end{proposition}

\begin{proof}
    Soit $(u_n) \in \ell^p(\N)^\N$ et $u_n = (u_k^n)_{k \in \N}$ une suite de Cauchy de $\ell^p(\N)$.

    Alors, $\forall \varepsilon > 0$, $\exists N \in \N$, $\forall m,n \geqslant N$, 
    $||u_n - u_m||_p < \varepsilon$.
    avec 
    
    $\displaystyle ||u_n - u_m||_p = \left( \sum_{k=0}^{+\infty} |u_k^n - u_k^m|^p \right)^{\frac{1}{p}}$.

    Soit $k \in \N$ fixé, alors la suite $(u_k^n)_{n \in \N}$ est de Cauchy dans $\R$, donc elle converge vers une limite $u_k \in \R$.

    On a alors $u = (u_k)_{k \in \N} \in \R^\N$.

    \uindent{Etape 2 } Montrons que $u \in \ell^p(\N)$.

    \svspace Soit $(u_n) \in \ell^p(\N)^\N$ une suite de Cauchy de $\ell^p(\N)$. Alors $(u_n)$ est bornée. 

    Donc, $\displaystyle \exists M > 0$, 
    $\forall n \in \N$, $\displaystyle \left( \sum_{k = 0 }^{+ \infty } |u_k^n|^p \right)^{\frac 1 p} \leqslant M$.

    Donc pour un $N \in \N$ fixé, on a $\displaystyle \left( \sum_{k=0}^{N} |u_k^n|^p \right)^{\frac 1 p} \leqslant M$.
    
    Donc $\displaystyle \sum_{k=0}^{N} |u_k|^p = \lim_{n \to +\infty} \sum_{k=0}^{N} |u_k^n|^p \leqslant M$

    Donc $\displaystyle \sum |u_k|^p$ converge, donc $u \in \ell^p(\N)$.

    \uindent{Etape 3 } Montrons que $u_n \to u$ dans $\ell^p(\N)$.

    \svspace Soit $\varepsilon > 0$, $N' \in \N$ tel que 
    $\displaystyle \left( \sum_{ k = 0 }^{N'} |u_k^n - u_k^m|^p \right)^{\frac 1 p} < \varepsilon$

    Et alors avec $m \to +\infty$, on a 
    $\displaystyle \left( \sum_{ k = 0 }^{N'} |u_k^n - u_k|^p \right)^{\frac 1 p} \leqslant \varepsilon$.

    Puis avec $N' \to +\infty$, on a 
    $\displaystyle \left( \sum_{ k = 0 }^{+\infty} |u_k^n - u_k|^p \right)^{\frac 1 p} \leqslant \varepsilon$.

    Donc $u_n \to u$ dans $\ell^p(\N)$.

\end{proof}

